\scanpage{121}
poids et la structure de polarisation sur $G^\circ$
\uline{se déterminent mutuellement}.

\deleted{Donc} [ À voir : \unsure{l'un|l'ens.|le gpe} des autom de $(G^\circ_{\mathbb R},\varepsilon_{\mathbb R},j_{\mathbb R})$
voire $(G^\circ_{\deleted{\unsure{\mathbb R}}},\varepsilon,\mu)$, voire $(G,\varepsilon,\mu)$ opère-t-il
transitivement sur $\operatorname{Pol}(G)$ ? \deleted{\unsure{En d'autres|Et de même}} \inserted{Par exemple}
\deleted{\unsure{termes}, y a-t-il} \inserted{est-il \unsure{impossible|possible}} \deleted{\unsure{un sens}} \deleted{\unsure{l'on}} de distinguer \deleted{\unsure{avec} \illegible{1}} \inserted{une}
polarisation ou un paquet de polarisations
préférentielles ? ]

Prenons p. ex. une catégorie de Motifs
géométriques sur $\mathbb F_q$, $q=p^n$, \deleted{\unsure{et}} Alors \deleted{\unsure{dans}}
il y a une classe de conjugaison d'hom
\[
\deleted{\unsure{j}}i\colon G_{m\,\mathbb C}\longrightarrow G_{\mathbb C}
\]
distinguée, avec
\[
\bar i\,i=j,\qquad \varepsilon i=\mathrm{id}
\]
% slanted
\margin{N.B. $i$ et \unsure{ses} conjugués sur $\mathbb Q$, \deleted{\unsure{ceux des} \illegible{2}} \inserted{\deleted{\unsure{division}}}, si $[K:\mathbb Q]=n$ \unsure{engendrent} $\operatorname{Hom}(G_{m\,\mathbb C},G^\circ_{\mathbb C})$}
définie, \deleted{\unsure{par}} si $G$ commutatif : $M\subset L$ \inserted{$\underset{K\ \mathrm{tot.}\ \text{réel}}{\cup}$ \unsure{quadrique|quadratique} tot. imag.},
par la forme $\lambda\mapsto\frac1n v_{\mathfrak p}(\lambda)$, où $\mathfrak p$ est un
idéal premier de $L$ sur $p$. $[\,v_{\mathfrak p}(\lambda)v_{\mathfrak p}(\bar\lambda)=
v_{\mathfrak p}(\lambda\bar\lambda)=v_{\mathfrak p}(q^{j^*(\lambda)})=j^*(\lambda);\,v_{\mathfrak p}(q)=1\,]$.
% slanted
\margin{\deleted{\unsure{est|cet}} \uuline{\unsure{canulé|annulé|Calculer}} $\longrightarrow$ cette forme \deleted{\unsure{originale}} n'est pas tjrs à valeurs entières, même sur le \unsure{s-groupe|ss-groupe} de $M$ où $j=0$, ex : prendre $q=p\ne2$ $L=\mathbb Q(\sqrt p)$, tel que nb de classes de $L$ soit $1$, p. ex $p:3,7,11,\ldots$}
En fait, il convient de ``normaliser'' $G_{\mathbb C}$ choix
d'un tel $i$, (NB. Un autom de $G^\circ$ \inserted{\fbox{sur $\mathbb Q$!}} qui \deleted{induit}
respecte $i$ est l'identité !). La structure d'effectivité
sur $M$ est déterminée grâce à $i$, celle
des $M_*$ (\unsure{le cat. de $h$|la cat. de $h$|$K$ ext. de $\mathbb Q$}) en particulier de $M_{\mathbb R},M_{\mathbb C}$,
par la connaissance du système de ses conjugués
% archive sheet number: 120
